\section{克莱因曲面基本群的表示}
\label{sec:02.3}

\subsection{实曲线基本群的线性表示}
\label{sec:02.3.1}

如 2.2 节所见，实曲线 $(X,\sigma)$ 的基本群是万有覆盖空间 $\tilde X$ 的自微分同胚群
\[
\pi^R_1((X,\sigma);x)=
\bigl\{f:\tilde X\longrightarrow\tilde X\ \bigm|\
\exists\sigma_f\in\{\Id_X;\sigma\},\ p\circ f=\sigma_f\circ p\bigr\}^{\mathrm{op}},
\]
其中 $f$ 覆盖 $\Id_X$ 或 $\sigma$。这个群的主要特征是：它把拓扑基本群含为正规子群，
该子群是到 $\{\Id_X,\sigma\}^{\mathrm{op}}\simeq\Z/2\Z$ 的一个自然满群同态的核：
\begin{equation}\label{eq:2.20}
1\longrightarrow\pi_1(X;x)\longrightarrow\pi^R_1((X,\sigma);x)
\longrightarrow\Z/2\Z\longrightarrow1.
\end{equation}
我们回忆：到 $\Z/2\Z$ 的满群同态称为 $\Z/2\Z$ 的一个增广；并且 $\tilde X$ 上有唯一的
黎曼面结构使拓扑覆盖映射 $p:\tilde X\to X$ 成为全纯映射。由于
$\sigma:X\to X$ 反全纯，满足 $p\circ f=\sigma\circ p$ 的映射
$f:\tilde X\to\tilde X$ 是反全纯的（当然，若 $p\circ f=p$ 则 $f$ 全纯），即
$\pi^R_1((X,\sigma);x)$ 通过或全纯或反全纯的变换作用在 $\tilde X$ 上，具体取决于
$f\in\pi^R_1((X,\sigma);x)$ 在 $X$ 上覆盖的是 $\Id_X$ 还是 $\sigma$。因此，在研究
$\pi^R_1((X,\sigma);x)$ 的复线性表示时，自然的要求是：$\pi^R_1((X,\sigma);x)$ 中
$\pi_1(X;x)\subset\pi^R_1((X,\sigma);x)$ 的元素通过复向量空间 $V$ 的 $\C$-线性变换
作用，而 $\pi^R_1((X,\sigma);x)\setminus\pi_1(X;x)$ 的元素通过 $V$ 的 $\C$-反线性变换
作用。这类变换特别地是 $\R$-线性的，并构成一个群，我们记作 $\GL^{\pm}(V)$；它把
$\GL(V)$ 含为指数 $2$ 的子群。确切地说，$\GL^{\pm}(V)$ 是 $V$ 的 $\C$-线性与
$\C$-反线性变换之群。上面所说「$\pi^R_1((X,\sigma);x)\setminus\pi_1(X;x)$ 的元素通过
$\GL^{\pm}(V)\setminus\GL(V)$ 中的变换作用在 $V$ 上，而 $\pi_1(X;x)$ 的元素通过
$\GL(V)$ 中的变换作用」，其确切含义是：存在群同态
$\rho:\pi^R_1((X,\sigma);x)\to\GL^{\pm}(V)$，使得下图可交换（即 $\rho$ 是两个增广之间的同态）：
\begin{equation}\label{eq:2.21}
\begin{tikzcd}[column sep=large, row sep=large]
1 \arrow[r] & \pi_1(X;x) \arrow[r] \arrow[d, "\rho\vert_{\pi_1(X;x)}"'] & \pi^R_1((X,\sigma);x) \arrow[r] \arrow[d, "\rho"'] & \Z/2\Z \arrow[r] \arrow[d, "="] & 1\\
1 \arrow[r] & \GL(V) \arrow[r] & \GL^{\pm}(V) \arrow[r] & \Z/2\Z \arrow[r] & 1
\end{tikzcd}
\end{equation}
其中映射 $\GL^{\pm}(V)\to\Z/2\Z$ 把 $V$ 的 $\C$-反线性变换送到 $\Z/2\Z$ 的非平凡元。
我们把这样的映射 $\rho:\pi^R_1((X,\sigma);x)\to\GL^{\pm}(V)$ 称为
$\pi^R_1((X,\sigma);x)$ 的一个\emph{线性表示}。若在 $V$ 上选取 Hermite 度量 $h$，我们
同样可以考察 $(V,h)$ 的 $\C$-线性与 $\C$-反线性等距变换之群 $\U^{\pm}(V,h)$，它把酉群
$\U(V,h)$ 含为指数 $2$ 的子群。于是按定义，\emph{酉表示}是增广之间的同态
\begin{equation}\label{eq:2.22}
\begin{tikzcd}[column sep=large, row sep=large]
1 \arrow[r] & \pi_1(X;x) \arrow[r] \arrow[d, "\rho\vert_{\pi_1(X;x)}"'] & \pi^R_1((X,\sigma);x) \arrow[r] \arrow[d, "\rho"'] & \Z/2\Z \arrow[r] \arrow[d, "="] & 1\\
1 \arrow[r] & \U(V,h) \arrow[r] & \U^{\pm}(V,h) \arrow[r] & \Z/2\Z \arrow[r] & 1
\end{tikzcd}
\end{equation}
复向量空间 $V$ 上的\emph{线性实结构}是 $V$ 的一个 $\C$-反线性对合 $\alpha$。特别地，
此时 $V$ 是 $V^\alpha$ 的复化，其中 $V^\alpha$ 是 $V$ 中由 $\alpha$ 的不动点构成的实
子空间：$V=V^\alpha\otimes_\R\C$，我们把对子 $(V,\alpha)$ 称为\emph{实向量空间}。复
向量空间 $V$ 上的线性实结构总是存在的（见习题 2.3.5）。若 $V$ 带有 Hermite 度量
$h$，则可假定 $\alpha$ 是 $h$ 的（反线性）等距，此时称 $\alpha$ 与 $h$ 相容。

一旦选定了实结构，群 $\GL^{\pm}(V)$ 与 $\U^{\pm}(V,h)$ 就具有典范的半直积结构。

\begin{lemma}\label{lem:2.3.1}
设 $(V,\alpha)$ 是带有相容 Hermite 度量的实向量空间。则 $V$ 的 $\C$-线性与
$\C$-反线性变换之群 $\GL^{\pm}(V)$ 同构于半直积
$\GL(V)\rtimes_\alpha\Z/2\Z$，其中 $\{\Id_V,\alpha\}\simeq\Z/2\Z$ 通过
$\Ad_\alpha u:=\alpha\circ u\circ\alpha^{-1}$ 作用在 $\GL(V)$ 上。同样，
$(V,h)$ 的 $\C$-线性与 $\C$-反线性等距之群 $\U^{\pm}(V,h)$ 同构于半直积
$\U(V)\rtimes_\alpha\Z/2\Z$。
\end{lemma}

注意 $\alpha$ 自然地也是 $\GL^{\pm}(V)$（相应地 $\U^{\pm}(V,h)$）中的元素，且在增广
映射 $\GL^{\pm}(V)\to\Z/2\Z$ 下被送到 $\alpha$；并且若 $\C$-线性变换 $u$ 在 $V$ 的
某个由 $\alpha$-不动点组成的基（实基）下的矩阵为 $A$，则
$\alpha\circ u\circ\alpha^{-1}$ 在该基下的矩阵为 $\bar A$。引理 2.3.1 的证明作为习题
2.3.6 留给读者。

\begin{definition}[线性表示与酉表示]\label{def:2.3.2}
$\pi^R_1((X,\sigma);x)$ 的\emph{线性}（相应地，\emph{酉}）表示，是从
$\pi^R_1((X,\sigma);x)$ 到 $\GL^{\pm}(V)$（相应地 $\U^{\pm}(V,h)$）的 $\Z/2\Z$ 的增广
之间的同态。确切地说，它是与增广同态 $\GL^{\pm}(V)\to\Z/2\Z$（相应地
$\U^{\pm}(V,h)\to\Z/2\Z$）在 (2.21)（相应地 (2.22)）的意义下相容的群同态
$\rho:\pi^R_1((X,\sigma);x)\to\GL^{\pm}(V)$（相应地
$\rho:\pi^R_1((X,\sigma);x)\to\U^{\pm}(V,h)$）。这类表示的集合记作
$\Hom_{\Z/2\Z}(\pi^R_1((X,\sigma);x);\GL^{\pm}(V))$（相应地
$\Hom_{\Z/2\Z}(\pi^R_1((X,\sigma);x);\U^{\pm}(V,h))$）。它带有自然的拓扑，该拓扑由
$\pi^R_1((X,\sigma);x)$ 到 $\GL^{\pm}(V)$ 的所有映射之空间的乘积拓扑诱导。
\end{definition}

注意：由于短正合列 (2.20)，$\pi^R_1((X,\sigma);x)$ 是有限生成的；而 $\GL^{\pm}(V)$ 是
有限维复向量空间 $V$ 的 $\R$-线性自同构群（后者带有典范拓扑）的子群，故定义 2.3.2 中
所说的拓扑，是有限多个 $\GL^{\pm}(V)$ 之积中的子空间拓扑。

接下来要问的自然问题是：两个这样的表示在什么时候应被认为是等价的？先验地，我们可以
用任何 $u\in\GL^{\pm}(V)$ 共轭给定的映射
$\rho:\pi^R_1((X,\sigma);x)\to\GL^{\pm}(V)$，而不改变 (2.21) 的可交换性，因为
$\GL(V)$ 在 $\GL^{\pm}(V)$ 中正规。然而，我们将考虑一个更严格的等价概念，其中要求
$u$ 落在 $\GL(V)$ 中。这将保证存在从后文要讨论的表示空间
$\Hom_{\Z/2\Z}(\pi^R_1((X,\sigma);x);\GL^{\pm}(V))/\GL(V)$ 到通常表示空间
$\Hom(\pi_1(X;x);\GL(V))/\GL(V)$ 的映射。事实上，它还蕴涵
$\pi^R_1((X,\sigma);x)$ 的表示空间只依赖于目标群（在我们的情形即 $\GL^{\pm}(V)$）的
$\Z/2\Z$-增广类，而不依赖于它作为 $\pi_1(X;x)$ 被 $\Z/2\Z$ 扩张的同构类。

在下文中，由于我们希望考虑具有 Hausdorff 拓扑的商空间，我们将把注意力限制在酉表示
空间上。上面的讨论直接适用于酉情形。

\begin{definition}[酉表示簇]\label{def:2.3.3}
$\pi^R_1((X,\sigma);x)$ 的两个酉表示称为\emph{等价}，若它们落在
$\U(V,h)$ 在 $\Hom_{\Z/2\Z}(\pi^R_1((X,\sigma);x);\U^{\pm}(V,h))$ 上的（连续）作用
$u\cdot\rho=\Ad_u\circ\rho$ 的同一条轨道之中。按定义，$\pi^R_1((X,\sigma);x)$ 到
$\U^{\pm}(V,h)$ 的\emph{酉表示簇}是轨道空间
\[
\Hom_{\Z/2\Z}(\pi^R_1((X,\sigma);x);\U^{\pm}(V,h))/\U(V,h),
\]
并赋予商拓扑。
\end{definition}

注意，尽管目标群是 $\U^{\pm}(V,h)$ 而非 $\U(V,h)$，我们仍把它称为
$\pi^R_1((X,\sigma);x)$ 的酉表示簇。为简化记号，我们一次性固定 $V$ 上的 Hermite 度量
$h$，并直接把群 $\U(V,h)$ 与 $\U^{\pm}(V,h)$ 记作 $\U(V)$ 与 $\U^{\pm}(V)$。由于我们对
$\pi^R_1((X,\sigma);x)$ 的表示之间等价关系所作的选择，存在从
$\pi^R_1((X,\sigma);x)$ 的酉表示簇到 $\pi_1(X;x)$ 的酉表示簇的映射：
\[
\begin{gathered}
\Phi:\ \Hom_{\Z/2\Z}(\pi^R_1((X,\sigma);x);\U^{\pm}(V))/\U(V)\\
\longrightarrow\Hom(\pi_1(X;x);\U(V))/\U(V),
\qquad [\rho]\longmapsto\bigl[\rho|_{\pi_1(X;x)}\bigr].
\end{gathered}
\]
在下一小节中，我们研究这个映射的一些性质。
